\section{Introduction}
\label{sec:intro}

Every timing-integrity monitor publishes numbers, and every monitor publishes
them on its own terms. A time-RAIM receiver reports a timing protection level
against a fault-free-plus-single-fault hypothesis set on a GNSS scenario of its
authors' choosing; a holdover analysis reports a coast bound against a clock model
and a horizon of its authors' choosing; an authenticated secure-time protocol
reports a cryptographic freshness guarantee against a threat model of its authors'
choosing. Each of these is a real, useful result. But there is no common ruler.
Two monitors that both claim an integrity risk of \num{e-6} may have been measured
against entirely disjoint fault sets, one of which happens to omit precisely the
fault the other survives; a monitor that claims to ``detect'' a delay attack and
one that claims to ``bound'' it may be scored as equals even though one of those
claims is provably impossible. What is missing is not another monitor. It is a
\emph{shared, open, versioned yardstick}---a fixed fault menu and a fixed scoring
rule---against which any monitor's \emph{stated} integrity risk can be compared to
its \emph{measured} undetected-error exceedance, on the same stimulus, with the
impossibility results of the field encoded so that no monitor can be credited for
an impossible feat. This paper defines that yardstick.

\paragraph{The gap, stated precisely.}
We are careful about the novelty, because the field is mature and almost every
ingredient here has a parent. The integrity taxonomy---nominal, unavailable,
misleading, hazardously misleading---is the Stanford--ESA integrity diagram for
GNSS position integrity \citep{tossaint2007}, whose alert-limit vocabulary comes
from the WAAS Minimum Operational Performance Standards \citep{rtcado229}; we cite
both as the scoring \emph{method} and specialize the diagram to the scalar time
error, claiming no new taxonomy. The integrity-risk-budgeting frame---compare a
\emph{stated} risk to a \emph{measured} exceedance---is the ARAIM lineage
\citep{blanch2015}; we inherit the framing, not a new estimator. Timing protection
levels themselves exist: Gioia's receiver-autonomous timing protection level
\citep{gioia2023} and the companion composed timing protection level
\citep{baweja2026cti}, of which the single-source conditional protection level
\citep{baweja2026tpl} is a special case, are monitors this benchmark can
\emph{score}, not results it restates. The delay-attack impossibility is
Mizrahi's \citep{mizrahi7384}, sharpened by the half-round-trip-time asymmetry
boundary of Narula and Humphreys \citep{narula2018}; we encode it, we do not
rediscover it. And the secure-time protocols---Network Time Security
\citep{rfc8915}, Roughtime \citep{roughtime}, and the fault-tolerant selection of
IEEE~802.1ASed \citep{ieee8021ased}---each carry a bespoke, per-protocol
evaluation, which is exactly the fragmentation the benchmark exists to make
comparable. What does not exist, \emph{to our knowledge}, is the open
cross-protocol conformance benchmark that sits above all of these as their common
ruler. We claim that superset object and nothing stronger: we do \emph{not} claim
the first secure-time protocol, the first time-RAIM, or that no timing protection
level exists---they all do, and the benchmark's value is precisely that it can
score them side by side.

\paragraph{Why a benchmark, of all things, makes a stable reference.}
There is a structural reason this object is worth building and worth citing. A
monitor makes an accuracy claim, and an accuracy claim can be wrong: it can be
optimistic against reality, and to trust it one must validate it against an
external truth. A benchmark, by contrast, makes \emph{no accuracy claim of its
own}. It does not assert that any clock is good, that any protection level is
tight, or that any monitor is safe; it asserts only that, on a fixed menu, a given
monitor's declared protection level did or did not bound its true error, and
tallies the outcome. Because it asserts nothing about accuracy, \emph{it cannot be
wrong about accuracy}. We call this property \emph{honesty-immunity}, and we take
it as the paper's organizing principle: a reference object that cannot overclaim
is a stable thing to cite, precisely because a citing author need
not take any accuracy claim on faith. The consequence, made explicit throughout,
is that this paper carries \emph{zero} \Validated{} tags---not because we failed to
validate something, but because there is nothing accuracy-shaped here to validate.
The one oracle we do run---an independent re-count of the region tallies---checks
the \emph{arithmetic of the scorer}, not the accuracy of any monitor, and is
labelled \Modelled{}~/~\code{InternalConsistency} accordingly.

\paragraph{The impossibility, made a type-level invariant.}
The benchmark's cleverest and most load-bearing feature is how it treats faults
that no monitor can detect. Mizrahi's result \citep{mizrahi7384} is that a
symmetric single-path delay is undetectable by any protocol, cryptographic or not;
the same holds for a replay inside the freshness window. A benchmark that scored
such a fault by asking ``did the monitor detect it?'' would be scoring an
impossible act, and would reward a monitor for falsely claiming to have detected
it. We refuse to represent that outcome at all. The scorecard's \code{Verdict}
enumeration has four values---\CovPass{}, \CovFail{}, \Absorbed{}, and
\Unabsorbed{}---and \emph{no} \textsf{Detected} value. For an undetectable
scenario, the only reachable verdicts are \Absorbed{} (the protection level was
large enough to carry the undetectable offset at every epoch) and \Unabsorbed{}
(it was not). It is therefore \emph{structurally impossible}, at the level of the
type system and machine-checked by the crate's tests, for the harness to report an
undetectable fault as detected. Mizrahi's impossibility is not merely cited in a
comment; it is compiled into the shape of the result. This turns a negative
result---``you cannot detect this''---into a positive scoring rule: a monitor is
rewarded for \emph{carrying enough protection level to absorb} an undetectable
offset, never for the impossible act of detecting it.

\paragraph{Contributions.}
The contribution is narrow and, we hope, durable.
\begin{enumerate}
  \item \textbf{The Stanford integrity diagram specialized to time}
  (\Cited{} method, \Proven{} specialization; Section~\ref{sec:method}). Each
  epoch is classified by $|\test-\utc|$ against the declared protection level and
  the application alert limit into \Nominal{}/\Unavail{}/\MI{}/\HMI{}, and a
  scenario's overbound-coverage verdict is driven by the \HMI{} rate alone---the
  hazardous case---while \MI{}, availability, and the nominal rate are reported as
  softer signals.
  \item \textbf{An open, parametric fault menu}
  (\Modelled{}, with a \Cited{} detectability partition;
  Section~\ref{sec:faultmenu}). Ten fault classes, each a true-error generator plus
  a detectability flag, with the undetectable set (symmetric delay, replay) fixed
  by Mizrahi and the asymmetric-delay offset fixed at half the path asymmetry by
  Narula and Humphreys.
  \item \textbf{The verdict-algebra impossibility encoding}
  (\Proven{}; Section~\ref{sec:verdict}). The four-value \code{Verdict} with no
  \textsf{Detected} variant makes reporting an undetectable fault as detected
  unrepresentable; the property is machine-checked.
  \item \textbf{The honesty-immunity argument}
  (Section~\ref{sec:immunity}). Why a benchmark that makes no accuracy claim is a
  stable reference cite, why this forces zero \Validated{} tags, and why the
  only oracle is an \code{InternalConsistency} re-count of the counting.
  \item \textbf{A demonstration that maps rather than flatters}
  (\Modelled{}; Section~\ref{sec:demo}). Scoring the companion composed protection
  level and a deliberately broken monitor against a nanosecond-scale menu, showing
  the benchmark simultaneously confirms the reference monitor is sound in its
  holdover regime \emph{and} reveals precisely where it needs adversarial
  extensions.
\end{enumerate}

\paragraph{Honesty scope.}
We separate the kinds of statement and label them at every occurrence. The
integrity taxonomy, the alert-limit vocabulary, the ARAIM budgeting frame, the
delay-attack impossibility, and every monitor and protocol we score are
\Cited{}. The verdict-algebra impossibility encoding is \Proven{}, machine-checked
by the crate's test suite. The scoring harness, the fault menu, and every scenario
magnitude are \Modelled{}: representative, not measured, carrying no certification
or technology-readiness implication of any kind. The word \Validated{} appears in
this paper only to be \emph{defined and then shown never to be earned}: a benchmark
makes no accuracy claim, so no row is validated, and the single oracle we run is an
\code{InternalConsistency} re-count that is explicitly \emph{not} \Validated{} and
\emph{not} an external accuracy dataset. There is no certified alert limit, no
DO-229/DO-178/ECSS conformance, no endorsement by any agency; the Stanford diagram
and the WAAS MOPS are cited as method sources only. Every number is regenerated
from the open engine by the commands of Section~\ref{sec:availability}.
